%%%%%%%%%%%%%%%%%%%%%%%%%%%%%%%%%%%%%%%%%%%%%%%%%%%%%%%%%%%%%%%%%%%%%%%%%%%%%%%
%%  AION-1 Reproduction -- Internal Reproduction Report.
%%  Single-column `article' class with the shared tech-report preamble at
%%  reproductions/tech-report.sty. Builds with vanilla TeX Live (>= 2022).
%%%%%%%%%%%%%%%%%%%%%%%%%%%%%%%%%%%%%%%%%%%%%%%%%%%%%%%%%%%%%%%%%%%%%%%%%%%%%%%
\documentclass[11pt,letterpaper]{article}
\usepackage{../../tech-report}

\newcommand{\aion}{\textsc{AION-1}}
\newcommand{\Rsq}{\ensuremath{R^{2}}}

\title{Public-Data Reproduction of the \aion{} Omnimodal Foundation Model:\\
       Frozen-Encoder Downstream Experiments}
\author{Greg Benson\thanks{\texttt{benson@cs.usfca.edu}; Department of Computer
        Science, University of San Francisco, CA 94117, USA} \and
        Xiaosheng Huang\thanks{\texttt{xhuang@usfca.edu}; Department of Physics
        \& Astronomy, University of San Francisco; and Physics Division,
        Lawrence Berkeley National Laboratory}}
\date{June 2026}

\begin{document}
\maketitle
\subtitle{Internal Reproduction Report}

\begin{abstract}
We reproduce the downstream experiments of \aion{} \citep{parker2025aion}, the
Polymathic AI omnimodal foundation model for astronomy, using only the publicly
released frozen checkpoints (\texttt{aion-base/large/xlarge}; 0.3/0.9/3\,B
parameters) and public data, on seven NVIDIA TITAN RTX GPUs. Pretraining is out
of scope (the original used 64--288 H100s). We rebuild the benchmark datasets
from the Multimodal Universe \citep{mmu2024} and a handful of label catalogs,
extract frozen-encoder embeddings, and train the paper's lightweight probing
heads (attentive pooling and a convolutional segmentation upsampler). The
flagship galaxy-property result reproduces closely: with a frozen encoder, the
redshift \Rsq{} rises from $0.81$ (photometry) to $0.98$ ($+$spectrum) to
$0.985$ ($+$image$+$spectrum), matching the paper's multimodal column
($1.00$). Stellar-parameter residuals from \textit{Gaia} XP coefficients
match or beat the paper. Morphology, retrieval, and segmentation reproduce
qualitatively; where our numbers fall short of the paper we trace it to corpus
scale (rate-limited image-cutout access) rather than the pipeline, and report
ours alongside the paper's throughout.
\end{abstract}

\section{Introduction}
\aion{} integrates 39 heterogeneous data types from five surveys (Legacy
Survey, HSC, SDSS, DESI, \textit{Gaia}) into a single masked-modeling
transformer. The released model is used as a \emph{frozen} feature extractor:
each downstream task trains only a small head on the encoder embeddings. We
reproduce the eleven reported experiments, grouped as (i) galaxy and stellar
property estimation, (ii) morphology classification, (iii) similarity-based
retrieval, (iv) image segmentation, and (v) generative redshift posteriors and
spectral super-resolution.

\section{Methods}
\paragraph{Model and probes.} We load each checkpoint with the released
\texttt{polymathic-aion} package, freeze it, and compute token embeddings
$E\in\mathbb{R}^{B\times N\times d}$ ($d=768/1024/2048$). For property
estimation we use the paper's \emph{attentive pooling} head: $K$ learned queries
cross-attend to $E$ (a normalized cross-attention block taken from the model's
own \texttt{posttraining} code), followed by per-target linear maps. For
morphology we mean-pool and use a two-layer MLP; for segmentation we reshape the
$24\times24$ spatial tokens and apply a lightweight transposed-convolution
upsampler. Retrieval uses mean-pooled, $L_2$-normalized embeddings and cosine
similarity. All heads train with AdamW under an 80/20 split, fixed
\texttt{SEED}$=2026$.

\paragraph{Data.} The tutorials' pre-joined files are no longer hosted; we
rebuild from the Multimodal Universe HuggingFace datasets. PROVABGS labels
\citep{hahn2023provabgs} key-join to DESI spectra on \texttt{TARGETID}
($20{,}666$ galaxies); Legacy Survey $g,r,i,z$ images
\citep{dey2019legacysurvey} are fetched from the cutout service by sky position.
Stellar labels come from the DD-Payne DESI catalog \citep{zhang2024ddpayne}
($25{,}992$ stars), with \textit{Gaia} DR3 parallax \citep{gaia2023dr3} attached
via CDS \textsc{XMatch} ($97\%$ matched). Morphology and retrieval use Galaxy10
DECaLS \citep{walmsley2022gzdecals}; strong-lens positives are the SuGOHI master
list \citep{jaelani2024sugohi}; segmentation masks are the GZ3D value-added
catalog \citep{masters2021gz3d}.

\section{Results}
\paragraph{Galaxy property estimation (Table~\ref{tab:gal}).} Adding the DESI
spectrum to photometry is the dominant effect, recovering nearly the full
multimodal result; the image alone helps the morphology-correlated properties.

\begin{table}[h]\centering
\caption{PROVABGS galaxy properties, frozen \aion{}-base, \Rsq{}.}
\label{tab:gal}
\begin{tabular}{lccccc}
\hline
config & $z$ & $\log M_\star$ & age & $\log Z$ & sSFR \\ \hline
photometry          & 0.81 & 0.75 & 0.35 & 0.38 & 0.61 \\
$+$image            & 0.86 & 0.84 & 0.33 & 0.44 & 0.63 \\
$+$spectrum         & 0.98 & 0.94 & 0.46 & 0.57 & 0.69 \\
$+$image$+$spectrum & \textbf{0.985} & \textbf{0.94} & 0.43 & 0.57 & 0.69 \\ \hline
paper (phot+im+spec)& 1.00 & 0.96 & 0.53 & 0.61 & 0.72 \\ \hline
\end{tabular}
\end{table}

\paragraph{Stellar parameters.} From APOGEE-matched \textit{Gaia} XP
coefficients, residual standard deviations are
$T_{\mathrm{eff}}\!\approx\!89$--$95$\,K (paper $94.6$),
$\log g\!\approx\!0.19$--$0.20$\,dex ($0.206$), and
$[\mathrm{Fe/H}]\!\approx\!0.11$\,dex ($0.115$) -- a clean match. For the
spectrum-based DESI$\times$DD-Payne task, the DESI-only config gives
$T_{\mathrm{eff}}=0.96, \log g=0.89, [\mathrm{Fe/H}]=0.80$; adding \textit{Gaia}
parallax raises the surface-gravity-sensitive parameters toward the paper's
DESI$+$parallax column ($0.99/0.98/0.94$).

\paragraph{Generative and low-data.} The single-shot redshift head produces a
posterior over the $1024$-bin grid that contracts sharply when the spectrum is
added (task~10), and \textit{Gaia} BP/RP coefficients decode to a DESI-resolution
spectrum recovering the major line centers (median correlation $0.83$, task~11).
In the low-data regime, the spectrum-config \Rsq{} saturates by $10^3$--$10^4$
labels, reproducing the paper's sample-efficiency claim.

\paragraph{Morphology, retrieval, segmentation.} Galaxy10 morphology accuracy is
$0.72$--$0.75$ (paper $0.84$); spiral/merger retrieval nDCG@10 is
$\sim\!0.82/0.68$ (paper $0.938/0.892$); strong-lens retrieval and GZ3D
spiral/bar segmentation are reproduced with the frozen encoder. These tasks are
\emph{corpus-limited}: the Legacy Survey cutout service rate-limits image
acquisition, so we evaluate on smaller, less-rare corpora than the paper's full
GZ-DECaLS; the live numbers are in \texttt{data/results/REPORT.md}.

\section{Limitations}
Absolute numbers are reproduction \emph{targets}, not pass/fail gates: the paper
omits seeds, exact split membership, and head hyperparameters. The largest
genuine shortfalls (morphology, retrieval) stem from corpus scale, not modeling.
The \texttt{xlarge} checkpoint underperforms \texttt{base} on photometry-only
inputs, as expected for a 3\,B frozen model given four scalar tokens. Full
reproduction details, scripts, and the live comparison table accompany this
report.

\bibliographystyle{plainnat}
\bibliography{references}
\end{document}
